% adopted from "Riddles in the Dark" template on Overleaf
% incorporated ideas from "Mailmerged Conference Name Cards" template 
\documentclass[grid,avery5371]{flashcards}
% Encoding
\usepackage[utf8]{inputenc}
\usepackage[T1]{fontenc}
% Fonts
\usepackage{ebgaramond}
% Layout
\usepackage{geometry}
\geometry{headheight=12pt,footskip=4pt}
\usepackage{multicol}
% Icons (modern version)
\usepackage{fontawesome5}
% URL formatting (load late)
\usepackage{hyperref}
% Headers
\usepackage{fancyhdr}
\pagestyle{fancy}
\fancyhf{}
\renewcommand{\headrulewidth}{0pt}
% specify deck variant
\newcommand{\deck}{floss}
\newcommand{\game}{emergence}
\chead{\small \game\ ({\sc \deck} deck)}
% Card styles
\cardbackstyle[\Huge]{plain}
\cardfrontstyle[\large]{headings}
% Metadata
\title{\game}
\author{Matthijs den Besten}
% read card info from CSV
\usepackage{datatool}
\DTLloaddb{cards}{cards.csv}
\begin{document}
\cardfrontfoot{\game\qquad\faCreativeCommons}
% ------------------------------------------------
% Game evaluation card
% ------------------------------------------------
\begin{flashcard}[\faMeh\quad Game Evaluation Criteria]{
\begin{itemize}\tiny \setlength{\itemsep}{.1ex}
\item Category Biases
\item Scale 1 to 7
\begin{multicols}{3}
\begin{itemize}
\item Clarity
\item Flow
\item Balance
\item Length
\item Integrity
\item Fun
\end{itemize}
\end{multicols}
\begin{multicols}{2}
\item Strongest Point
\item Weakest Point
\item One Change
\item Comparable Games
\end{multicols}
\end{itemize}
}
{\bfseries\small \game --- Game Evaluation}
\small
\begin{tabular}{l|c|c|c|c|}
 & 1 & 3 & 5 & 7\\
\hline
{\em Clarity} & Opaque & Muddy & Transparent & Water-clear \\
{\em Flow} & Cumbrous & Fraught & Smooth & Natural \\
{\em Balance} & Broken & Fluky & Sensible & Fair\\
{\em Length} & Wrong & Unfit & Apt & Perfect \\
{\em Integrity} & Eristic & Erratic & Coherent & Sound \\
{\em Fun} & Offputting & Boring & Engaging & Exciting \\
\hline
\end{tabular}
\end{flashcard}
% ------------------------------------------------
% Generate cards from CSV
% ------------------------------------------------
\DTLforeach{cards}{
\Description=description,
\Category=family,
\Id=period_id,
\Tool=discovery,
\Dependson=dependson
}
{
\begin{flashcard}[\deck:\quad\Id\quad\DTLifnull{\Tool}{}{\Tool}]{
\Description
Reqs: \Dependson
}
\faTasks\quad{\sc \Tool}
\end{flashcard}
}
\end{document}
